\chapter{Strategi AI untuk Bisnis}
\label{chap:business}

Mengadopsi AI bukan sekadar membeli lisensi ChatGPT Enterprise. Ini membutuhkan perubahan budaya kerja.

\section{Framework Adopsi AI}

\begin{enumerate}
    \item \textbf{Educate:} Latih karyawan tentang literasi AI dasar (seperti membaca buku ini).
    \item \textbf{Experiment:} Izinkan tim kecil membuat PoC (Proof of Concept) risiko rendah.
    \item \textbf{Expand:} Skalakan eksperimen yang berhasil ke seluruh departemen.
    \item \textbf{Embed:} Integrasikan AI ke dalam inti produk bisnis.
\end{enumerate}

\section{Menghitung ROI AI}

Investasi AI mahal. Bagaimana menghitung untungnya?

$$ \text{ROI} = \frac{\text{Nilai Waktu yang Dihemat} + \text{Pendapatan Baru}}{\text{Biaya API} + \text{Biaya Tim}} $$

\section{Studi Kasus Sektoral}

\subsection{Kesehatan}
Dokter menggunakan AI untuk menulis rangkuman medis otomatis, menghemat 2 jam per hari untuk administrasi.

\subsection{Keuangan}
Bank menggunakan AI untuk mendeteksi transaksi kartu kredit mencurigakan dalam milidetik.

\subsection{E-Commerce}
Rekomendasi produk yang dipersonalisasi meningkatkan penjualan hingga 30\%.

\section{Risiko Shadow AI}

Shadow AI adalah ketika karyawan menggunakan tools AI gratisan tanpa izin IT.
\textbf{Bahaya:} Data perusahaan bocor ke publik.
\textbf{Solusi:} Sediakan tool resmi yang aman agar karyawan tidak mencari alternatif liar.
